\begin{thebibliography}{99}
\bibitem[{Abaee(2026a)}]{Abaee2026a} Abaee, A. (2026a). \emph{Governance delay and the stability of harvested stocks: mobilising and protective feedback rules, and the review interval as a design parameter}. Working paper (v41). \url{https://github.com/MIKEAA2020/general-sustainability/blob/2efaa62083f2007dec3446c15c475af88f0925f3/arena%20agent%201/paper%20rewrites/latex/paper4_delay_dynamics_v41.pdf}.
\bibitem[{Abaee(2026b)}]{Abaee2026b} Abaee, A. (2026b). \emph{Aggregation is a claim, not a presentation: the geometry of the acceptance gap between compensatory and noncompensatory transition assessment}. Working paper (v71). \url{https://github.com/MIKEAA2020/general-sustainability/blob/2efaa62083f2007dec3446c15c475af88f0925f3/arena%20agent%201/paper%20rewrites/latex/paper06_assessment_separation_v71.pdf}.
\bibitem[{Abaee(2026c)}]{Abaee2026c} Abaee, A. (2026c). \emph{Typed Flux Ledgers and Depletion Arithmetic: Conservation, Componentwise Diagnostics, and the Semantics of Depletion Horizons}. Working paper. PDF pinned to GitHub commit \url{https://github.com/MIKEAA2020/general-sustainability/blob/2efaa62083f2007dec3446c15c475af88f0925f3/paper3/latest-2026-09-19/manuscript/paper3_material_ledgers_v50.pdf}.
\bibitem[{Anderies(2000)}]{Anderies2000} Anderies, J. M. (2000). On modeling human behavior and institutions in simple ecological economic systems. \emph{Ecological Economics}, 35(3), 393--412.
\bibitem[{Becker et al.(2017)}]{Becker2017} Becker, W., Saisana, M., Paruolo, P. \& Vandecasteele, I. (2017). Weights and importance in composite indicators: Closing the gap. \emph{Ecological Indicators}, 80, 12--22.
\bibitem[{Blomqvist et al.(2013)}]{Blomqvist2013} Blomqvist, L., Brook, B. W., Ellis, E. C., Kareiva, P. M., Nordhaus, T. \& Shellenberger, M. (2013). Does the shoe fit? Real versus imagined ecological footprints. \emph{PLoS Biology}, 11(11), e1001700.
\bibitem[{Borucke et al.(2013)}]{Borucke2013} Borucke, M., Moore, D., Cranston, G., Gracey, K., Iha, K., Larson, J., Lazarus, E., Morales, J. C., Wackernagel, M. \& Galli, A. (2013). Accounting for demand and supply of the biosphere's regenerative capacity: The National Footprint Accounts' underlying methodology and framework. \emph{Ecological Indicators}, 24, 518--533.
\bibitem[{Brander and Taylor(1998)}]{BranderTaylor1998} Brander, J. A. \& Taylor, M. S. (1998). The simple economics of Easter Island: A Ricardo--Malthus model of renewable resource use. \emph{American Economic Review}, 88(1), 119--138.
\bibitem[{Calvin and Bond-Lamberty(2018)}]{CalvinBondLamberty2018} Calvin, K. \& Bond-Lamberty, B. (2018). Integrated human--Earth system modeling --- state of the science and future directions. \emph{Environmental Research Letters}, 13(6), 063006.
\bibitem[{Cohen(1995)}]{Cohen1995} Cohen, J. E. (1995). Population growth and Earth's human carrying capacity. \emph{Science}, 269(5222), 341--346.
\bibitem[{Dasgupta(2021)}]{Dasgupta2021} Dasgupta, P. (2021). \emph{The Economics of Biodiversity: The Dasgupta Review.} London: HM Treasury.
\bibitem[{Donges et al.(2017)}]{Donges2017} Donges, J. F., Winkelmann, R., Cornell, S. E., Lucht, W., Dyke, J. G., Rockstr\"om, J., Heitzig, J. \& Schellnhuber, H. J. (2017). Closing the loop: Reconnecting human dynamics to Earth system science. \emph{The Anthropocene Review}, 4(2), 151--157.
\bibitem[{Erb et al.(2017)}]{Erb2017} Erb, K.-H., Luyssaert, S., Meyfroidt, P., Pongratz, J., Don, A., Kloster, S., Kuemmerle, T., Fetzel, T., Fuchs, R., Hasegawa, T., Haberl, H., et al. (2017). Land management: data availability and process understanding for global change studies. \emph{Global Change Biology}, 23(2), 512--533.
\bibitem[{Fischer et al.(2022)}]{Fischer2022} Fischer, S. M., Joy, M. K., Abrahamse, W., Milfont, T. L. \& Petherick, L. M. (2022). The use and misuse of composite environmental indices. \emph{bioRxiv}, 2022.03.15.484501. \url{https://doi.org/10.1101/2022.03.15.484501}.
\bibitem[{Global Footprint Network(2009)}]{GFNStandards2009} Global Footprint Network (2009). \emph{Ecological Footprint Standards 2009}. Global Footprint Network. \url{https://www.footprintnetwork.org/content/uploads/2019/05/Ecological_Footprint_Standards_2009.pdf}.
\bibitem[{Global Footprint Network(2019)}]{GFNGuidebook2019} Global Footprint Network (2019). \emph{Working Guidebook to the National Footprint Accounts 2019}. Global Footprint Network. \url{https://www.footprintnetwork.org/content/uploads/2019/05/National_Footprint_Accounts_Guidebook_2019.pdf}.
\bibitem[{Galli et al.(2016)}]{Galli2016} Galli, A., Giampietro, M., Goldfinger, S., Lazarus, E., Lin, D., Saltelli, A., Wackernagel, M. \& M\"uller, F. (2016). Questioning the ecological footprint. \emph{Ecological Indicators}, 69, 224--232.
\bibitem[{Giampietro and Saltelli(2014)}]{GiampietroSaltelli2014} Giampietro, M. \& Saltelli, A. (2014). Footprints to nowhere. \emph{Ecological Indicators}, 46, 610--621.
\bibitem[{Haberl et al.(2007)}]{Haberl2007} Haberl, H., Erb, K.-H., Krausmann, F., Gaube, V., Bondeau, A., Plutzar, C., Gingrich, S., Lucht, W. \& Fischer-Kowalski, M. (2007). Quantifying and mapping the human appropriation of net primary production in earth's terrestrial ecosystems. \emph{PNAS}, 104(31), 12942--12947.
\bibitem[{Henderson and Loreau(2018)}]{HendersonLoreau2018} Henderson, K. \& Loreau, M. (2018). How ecological feedbacks between human population and land cover influence sustainability. \emph{PLoS Computational Biology}, 14(8), e1006389.
\bibitem[{Hutchings and Reynolds(2004)}]{HutchingsReynolds2004} Hutchings, J. A. \& Reynolds, J. D. (2004). Marine fish population collapses: consequences for recovery and extinction risk. \emph{BioScience}, 54(4), 297--309.
\bibitem[{Hutchinson(1948)}]{Hutchinson1948} Hutchinson, G. E. (1948). Circular causal systems in ecology. \emph{Annals of the New York Academy of Sciences}, 50(4), 221--246.
\bibitem[{Krausmann et al.(2013)}]{Krausmann2013} Krausmann, F., Erb, K.-H., Gingrich, S., Haberl, H., Bondeau, A., Gaube, V., Lauk, C., Plutzar, C. \& Searchinger, T. D. (2013). Global human appropriation of net primary production doubled in the 20th century. \emph{PNAS}, 110(25), 10324--10329.
\bibitem[{Lo et al.(2026)}]{Lo2026NFA} Lo, K., Basnet, N., Miller, E., Srikanthalingam, B., Foley, B., Van Berkum, J. L., Long, A. H., Toneva, P., Ermina, M. \& Klinkenberg, C. (2026). \emph{National Ecological Footprint and Biocapacity Accounts, Edition 2026 (1961--2025): Results and metadata, version 1.0.} Prepared for Footprint Data Foundation by researchers at York University and the University of Iceland. Workbook: \url{https://footprint.info.yorku.ca/files/2026/04/NEFBA_Data_global_national_ecological_footprint_biocapacity.xlsx}; release notes: \url{https://footprint.info.yorku.ca/files/2026/04/NEFBA_ReleaseNotes.pdf}; data page: \url{https://footprint.info.yorku.ca/data/}.
\bibitem[{May(1973)}]{May1973} May, R. M. (1973). \emph{Stability and Complexity in Model Ecosystems.} Princeton: Princeton University Press.
\bibitem[{Meadows et al.(1972)}]{Meadows1972} Meadows, D. H., Meadows, D. L., Randers, J. \& Behrens, W. W. (1972). \emph{The Limits to Growth.} New York: Universe Books.
\bibitem[{Meadows et al.(1974)}]{Meadows1974} Meadows, D. L., Behrens, W. W., Meadows, D. H., Naill, R. F., Randers, J. \& Zahn, E. (1974). \emph{Dynamics of Growth in a Finite World.} Cambridge, MA: Wright-Allen Press.
\bibitem[{Meadows et al.(2004)}]{Meadows2004} Meadows, D. H., Randers, J. \& Meadows, D. L. (2004). \emph{The Limits to Growth: The 30-Year Update.} White River Junction, VT: Chelsea Green.
\bibitem[{Meyer(1994)}]{Meyer1994} Meyer, P. S. (1994). Bi-logistic growth. \emph{Technological Forecasting and Social Change}, 47(1), 89--102.
\bibitem[{Meyer and Ausubel(1999)}]{MeyerAusubel1999} Meyer, P. S. \& Ausubel, J. H. (1999). Carrying capacity: A model with logistically varying limits. \emph{Technological Forecasting and Social Change}, 61(3), 209--214.
\bibitem[{Motesharrei et al.(2016)}]{Motesharrei2016} Motesharrei, S., Rivas, J., Kalnay, E., Asrar, G. R., Busalacchi, A. J., Cahalan, R. F., Cane, M. A., Colwell, R. R., Feng, K., Franklin, R. S., et al. (2016). Modeling sustainability: Population, inequality, consumption, and bidirectional coupling of the Earth and human systems. \emph{National Science Review}, 3(4), 470--494.
\bibitem[{Nardo et al.(2008)}]{Nardo2008} Nardo, M., Saisana, M., Saltelli, A., Tarantola, S., Hoffmann, A. \& Giovannini, E. (2008). \emph{Handbook on Constructing Composite Indicators: Methodology and User Guide.} Paris: OECD Publishing.
\bibitem[{Neubauer et al.(2013)}]{Neubauer2013} Neubauer, P., Jensen, O. P., Hutchings, J. A. \& Baum, J. K. (2013). Resilience and recovery of overexploited marine populations. \emph{Science}, 340(6130), 347--349.
\bibitem[{Our~World~in~Data(2026a)}]{OWIDCereal2026} Our~World~in~Data (2026a). \emph{Cereal yield.} FAOSTAT data accessed through Our World in Data; grapher page last updated 25 February 2026; accessed 24 October 2026. \url{https://ourworldindata.org/grapher/cereal-yield}.
\bibitem[{Our~World~in~Data(2026b)}]{OWIDCropland2026} Our~World~in~Data (2026b). \emph{Cropland area.} FAOSTAT data accessed through Our World in Data; grapher page last updated 25 February 2026; accessed 24 October 2026. \url{https://ourworldindata.org/grapher/cropland-area}.
\bibitem[{Pearl(1925)}]{Pearl1925} Pearl, R. (1925). \emph{The Biology of Population Growth.} New York: Alfred A. Knopf.
\bibitem[{Poeplau et al.(2011)}]{Poeplau2011} Poeplau, C., et al. (2011). Temporal dynamics of soil organic carbon after land-use change in the temperate zone--carbon response functions as a model approach. \emph{Global Change Biology}, 17(7), 2415--2427.
\bibitem[{Poorter et al.(2016)}]{Poorter2016} Poorter, L., et al. (2016). Biomass resilience of Neotropical secondary forests. \emph{Nature}, 530(7590), 211--214.
\bibitem[{Schaefer(1954)}]{Schaefer1954} Schaefer, M. B. (1954). Some aspects of the dynamics of populations important to the management of the commercial marine fisheries. \emph{Bulletin of the Inter-American Tropical Tuna Commission}, 1(2), 25--56.
\bibitem[{Scheffer et al.(2001)}]{Scheffer2001} Scheffer, M., Carpenter, S., Foley, J. A., Folke, C. \& Walker, B. (2001). Catastrophic shifts in ecosystems. \emph{Nature}, 413(6856), 591--596.
\bibitem[{Simon(1996)}]{Simon1996} Simon, J. L. (1996). \emph{The Ultimate Resource 2.} Princeton: Princeton University Press.
\bibitem[{United Nations et al.(2021)}]{UN2021SEEA_EA} United Nations et al. (2021). \emph{System of Environmental-Economic Accounting---Ecosystem Accounting (SEEA EA).} White cover publication, pre-edited text subject to official editing. New York: United Nations. \url{https://seea.un.org/ecosystem-accounting}.
\bibitem[{United Nations(2022)}]{UN2022} United Nations (2022). \emph{World Population Prospects 2022: Summary of Results.} New York: Department of Economic and Social Affairs, Population Division.
\bibitem[{van den Bergh and Grazi(2015)}]{vdBerghGrazi2015} van den Bergh, J. C. J. M. \& Grazi, F. (2015). Reply to the first systematic response by the Global Footprint Network to criticism: a real debate finally? \emph{Ecological Indicators}, 58, 458--463.
\bibitem[{van den Bergh and Rietveld(2004)}]{vdBerghRietveld2004} van den Bergh, J. C. J. M. \& Rietveld, P. (2004). Reconsidering the limits to world population: Meta-analysis and meta-prediction. \emph{BioScience}, 54(3), 195--204.
\bibitem[{Wackernagel and Rees(1996)}]{WackernagelRees1996} Wackernagel, M. \& Rees, W. (1996). \emph{Our Ecological Footprint: Reducing Human Impact on the Earth.} New Society Publishers.
\bibitem[{Wilson(2016)}]{Wilson2016} Wilson, E. O. (2016). \emph{Half-Earth: Our Planet's Fight for Life.} New York: Liveright.
\end{thebibliography}

\end{document}
